\documentclass[11pt]{article}
\usepackage[margin=1.1in]{geometry}
\usepackage{amsmath,amssymb}
\usepackage{booktabs}
\usepackage{hyperref}
\newcommand{\D}{\mathsf{D}}
\newcommand{\prox}{\operatorname{prox}}
\newcommand{\freg}{f_{\mathrm{reg}}}
\newcommand{\code}[1]{\texttt{#1}}
\title{Recovering the implicit regularizer of PIRATE+:\\ a self-contained task description}
\author{G. Provencher Langlois}
\date{August 26, 2026}
\begin{document}
\maketitle

\section{What we already know}

Plug-and-play methods insert a trained denoiser $\D$ into an iterative reconstruction, using it in place of the gradient step of a regularizer. PIRATE (Hu, Gan, Sun, An, Kamilov, ICLR 2024) does this for deformable image registration: it minimizes $g + h$ (their Eq.~5), and its iteration (their Eq.~6) takes the step $\phi - \gamma(\nabla g + \alpha\nabla r + \tau(\phi - \D(\phi)))$, so the implicit regularizer is the term whose gradient is the residual $\phi - \D(\phi)$. The paper justifies this reading through the identity for the exact minimum-MSE denoiser (their Eq.~11), $\phi - \D(\phi) = -\sigma^2\nabla\log p_z(\phi)$, which does make the residual a gradient. But that identity is a statement about the ideal denoiser. Their $\D$ is a trained network, and PIRATE+ is fine-tuned on a registration loss and so is the MMSE denoiser of no prior, so the identity has no referent for it. Whether the trained residual is a gradient of anything is neither stated nor tested.

We tested it. Two facts stand from the earlier experiments (recorded in \code{numerics/pnp\_reg/}, summarized in \code{\_\_notes/note\_experiments\_summary.tex}).

First, the released AWGN denoiser fails: its Jacobian is far from symmetric (asymmetry ratio $\rho = 0.44$), so no function has the residual as its gradient, and the objective PIRATE writes down does not exist for PIRATE's own denoiser.

Second, the fine-tuned variant PIRATE+ (the same network after deep-equilibrium fine-tuning) nearly passes: $\rho = 0.016$, symmetric part positive definite, spectrum in $[0.95, 1.05]$. On its operating distribution PIRATE+ is close to affine (a fixed linear map plus a constant) to $0.2\%$. An affine denoiser has a \emph{quadratic} implicit regularizer, which we then wrote in closed form from the measured linear part $A$ (with symmetric part $S$) and constant $b$:
\begin{equation}
\nabla\freg(y) = S^{-1}(y - b) - y .
\label{eq:quad}
\end{equation}
This closed form passed a held-out prediction (prox-identity residual $0.597$ against a predicted ceiling $0.704$). The relevant curvatures sit in $[-0.05, +0.06]$: small, and slightly negative in some directions.

\section{The experiment}

We recover $\freg$ a second way, by a fit that assumes nothing about its form, and check that it returns the same quadratic. The fit uses only forward passes of PIRATE+.

The identity is one line. If $\D = \prox_{\freg}$, then $y = \D(z)$ satisfies the optimality condition $z - y = \nabla\freg(y)$, i.e.
\begin{equation}
z = \nabla G(y), \qquad G := \freg + \tfrac12\|\cdot\|^2 ,
\label{eq:grad}
\end{equation}
and $G$ is convex whenever $\freg + \tfrac12\|\cdot\|^2$ is (which holds here, since $S \succ 0$). So the denoiser hands us gradient data for a convex potential: sample fields $z_k$, set $y_k = \D(z_k)$, and fit a convex network $G_\theta$ so that $\nabla G_\theta(y_k) = z_k$ in least squares. The regularizer is read off as $f_\theta = G_\theta - \tfrac12\|\cdot\|^2$.

The point of the experiment is what $f_\theta$ turns out to be. If PIRATE+ is affine, then $f_\theta$ should come out essentially quadratic: near-constant curvature, and a gradient field that matches \eqref{eq:quad}. The network is free to bend; the claim is that it does not.

\section{Why it matters}

Three reasons, in order.

The recovery is our headline method, and this is the first time we run it on a released operator rather than a hand-built one. It has to work here.

It confirms the affine finding by an independent route. Equation~\eqref{eq:quad} came from inverting a measured linear map; $f_\theta$ comes from a network that could have represented any nonconvex regularizer and was told nothing about quadratics. If the two agree, the affine account is not an artifact of assuming affinity.

It produces the explicit object PIRATE+ uses only implicitly: a function on deformation fields that we can evaluate, differentiate, and plot. With it, the PIRATE+ update becomes descent on a written-down objective, which is the statement the PIRATE paper assumes and cannot make.

State the result plainly in the writeup: because the signal is quadratic, this is a confirmation, not a discovery. The discovery was the affinity itself. This experiment closes it.

\section{Prerequisites}

\begin{itemize}
\item Repository \code{numerics/}. Released PIRATE weights are vendored at \code{numerics/ext/pirate/} (upstream \code{wustl-cig/PIRATE-code}, MIT). PIRATE+ is the fine-tuned checkpoint there; loading is shown in \code{ext/pirate/inference\_PIRATEplus.py}.
\item Python environment \code{lpn\_env} (\code{\textasciitilde/miniforge3/envs/lpn\_env/bin/python}); PyTorch, CPU is fine.
\item Existing code to reuse, not rewrite:
  \begin{itemize}
  \item \code{numerics/tv\_pm/tvpm/}: the gradient-supervised training path already used for the total-variation example --- \code{recover.build\_model} (a convolutional input-convex network), the training loop in \code{recover.run}, \code{dataset} for patching. Import these; do not edit \code{tv\_pm/} (it is a finished experiment).
  \item \code{numerics/pnp\_reg/pnpreg/readout.py}: \code{value(G,y)} and \code{grad(G,y)} implement $f_\theta = G_\theta - \tfrac12\|\cdot\|^2$ and its gradient.
  \item \code{numerics/pnp\_reg/pnpreg/affine.py}: the closed-form quadratic \eqref{eq:quad}, its symmetric part $S$, and the operating-distribution sampler \code{operating\_inputs}. This is the reference the fit is compared against; numbers are in \code{results/affine\_metrics.json}.
  \end{itemize}
\item New code goes in \code{numerics/pnp\_reg/pnpreg/} with a distinct checkpoint prefix; nothing under \code{tv\_pm/} or \code{ext/} is written.
\end{itemize}

\section{Data and weights (already in the repo)}

Nothing needs to be downloaded. The PIRATE release is vendored verbatim at \code{numerics/ext/pirate/} (a shallow clone of \code{wustl-cig/PIRATE-code}, commit \code{4901b71}, MIT license; provenance in \code{numerics/ext/PROVENANCE.md}). Do not edit anything under \code{ext/}.

Two denoiser checkpoints ship, both about $5$\,MB:
\begin{itemize}
\item PIRATE (the AWGN denoiser): \code{ext/pirate/pretrained\_model/AWGN\_denoiser/OASIS.pth.tar}.
\item PIRATE+ (the deep-equilibrium fine-tuned denoiser): \code{ext/pirate/pretrained\_model/PIRATEplus/OASIS.pth.tar}. It is the \emph{same} twelve DnCNN tensors as PIRATE, stored under a \code{PIRATE.dnn.} key prefix; the two differ only by the fine-tuning, so the comparison between them is exactly matched.
\end{itemize}

The example registration field is at \code{ext/pirate/data/field.h5py} (shape $(1,3,80,96,112)$, about $10$\,MB); \code{fixed.nii.gz} and \code{moving.nii.gz} in the same folder are the image pair, needed only for the optional plug-back check. The DnCNN class is \code{ext/pirate/model/base.py}.

Do not reimplement the loading. Our package already wraps it: \code{numerics/pnp\_reg/pnpreg/paths.py} holds the three paths, and \code{pnpreg/probe\_targets.py} provides \code{load\_pirate\_dncnn(which)} with \code{which} in \code{\{"pirate","pirate\_plus"\}} (returns the frozen network) and \code{load\_field(crop=None)}. The denoiser is $\D(z) = z - \mathrm{dnn}(z)$, since the network predicts the noise. There is one checkpoint per denoiser, trained at their $\sigma = 1$, where the noise standard deviation is $1$ (their noise is $\sigma^2\,\mathrm{randn}$); this is the single operating point, and it is the one \code{affine.operating\_inputs} samples.

\section{Procedure}

\begin{enumerate}
\item \textbf{Data.} Draw fields $z_k$ from PIRATE+'s operating distribution: the released registration field plus unit-variance Gaussian noise (the sampler is \code{affine.operating\_inputs}; the noise level is their $\sigma = 1$). Set $y_k = \D(z_k)$ by forward passes of PIRATE+. Because the regularizer acts patchwise, cut $(z_k, y_k)$ into patches, as the total-variation experiment does; a few fields give hundreds of thousands of patch pairs. Hold out a fresh set of fields for evaluation.
\item \textbf{Fit.} Train the convex network $G_\theta$ so that $\nabla G_\theta(y_k) = z_k$ in least squares, reusing \code{recover.build\_model} and \code{recover.run}. Use the same budget as the total-variation run. Keep the best-validation checkpoint.
\item \textbf{Read off} $f_\theta = G_\theta - \tfrac12\|\cdot\|^2$ via \code{readout}.
\item \textbf{Compare to the quadratic.} On held-out operating inputs, compare $\nabla f_\theta(y)$ against the closed form \eqref{eq:quad} from \code{affine.py}: relative $L^2$ of the two gradient fields, and the curvature spectrum of $f_\theta$ (curvatures should be near-constant and lie in the measured $[-0.05, +0.06]$ band, not spread out).
\item \textbf{Certificates without ground truth.} Report the held-out proximal residual $\|\nabla G_\theta(\D(z)) - z\|$ (median and a high quantile --- the median hides localized failures). Re-probe $\prox_{f_\theta} = \nabla G_\theta^{\,*}$ with the Experiment~2 machinery (\code{pnpreg/probe.py}): by construction it must read asymmetry $\rho = 0$, and its symmetric-part spectrum should reproduce $S$.
\end{enumerate}

\section{What to expect}

\begin{itemize}
\item $f_\theta$ is quadratic to within the fit error: gradient field matching \eqref{eq:quad} at the percent level, curvature spectrum tight around the measured band.
\item Held-out residual near the percent scale set by $\rho = 0.016$. As a control, the same fit on the AWGN denoiser ($\rho = 0.44$) should be worse by roughly $30\times$; if the residuals track the two asymmetry ratios, the probe and the recovery corroborate each other.
\item If instead $f_\theta$ shows structure beyond quadratic above $0.2\%$, that contradicts Experiment~3.1 and must be chased before anything is written --- start from whether the fit or the affinity measurement is at fault.
\end{itemize}

\section{Deliverables}

\begin{enumerate}
\item The function $f_\theta$: value and curvature along field slices, and its comparison against \eqref{eq:quad}, as one figure.
\item A short table: gradient-field relative $L^2$ (fit vs.\ closed form), curvature range, held-out residual (median and high quantile) for PIRATE+ and for the AWGN control.
\item Optional and only if cheap: the plug-back check --- replace $\D$ by $\prox_{f_\theta}$ inside PIRATE's own iteration on their example data and confirm registration accuracy is unchanged. Success gives a PIRATE variant with the same performance whose update provably descends the explicit $f_\theta$.
\end{enumerate}

\section{Pitfalls}

\begin{itemize}
\item Restrict every error to the sampled range of $\D$. The regularizer is pinned only where the denoiser was probed; values off that range are inductive bias, not measurement.
\item Report the high quantile of the residual, not only the median. In the earlier runs the median sat near $4.5\%$ while the $95$th percentile was $16.5\%$, and the tail carried the real information.
\item Do not read the recovered curvature from pointwise second derivatives of the network; on a smooth-activation network they carry spatial wiggle that swamps small curvatures. Fit slopes to gradient components over the operating range instead, as in \code{note\_experiments\_summary.tex}, Section~5 (the ``right instrument'' paragraph of Experiment~3.2).
\item Keep the noise convention theirs: their $\sigma = 1$ means noise standard deviation $1$, about half the field scale. This is the operating point; do not rescale it.
\end{itemize}

\end{document}
